% -*- coding: utf-8 -*-
% !TEX program = xelatex
\documentclass[plain,noanswer,medmath]{../../template/jnuexam}
\usepackage{../../template/kysx}

\begin{document}


\examtitle{2000}{4}


\loadproblems{3}{2000P1}%载入试卷一
\loadproblems{3}{2000P3}%载入试卷三


\makepart{填空题}{本题共5小题，每小题3分，满分15分}


\begin{problem}
$\int \frac{\arcsin \sqrt{x}}{\sqrt{x}} \dx =$ \fillin{}．
\end{problem}


\begin{problem}
若$a > 0$, $b > 0$均为常数，则$\lim_{x \to 0} \left(\frac{a^x + b^x}{2} \right)^{\frac{3}{x}} =$ \fillin{}．
\end{problem}


\begin{problem}
设$\alpha = (1,0, - 1)^T$，矩阵$A = \alpha \alpha^T$，$n$为正整数，则$|kE - A^n| =$ \fillin{}．
\end{problem}


%【类似数学三第一(3)题】
\begin{problem}
已知$4$阶矩阵$A$相似于$B$，$A$的特征值为$2,3,4,5$，$E$为$4$阶单位矩阵，则行列式$|B - E| =$ \fillin{}．
\end{problem}


\useproblem{3}{1}{5}%【同试卷三第一(5)题】


\makepart{选择题}{本题共5小题，每小题3分，共15分}


\useproblem{3}{2}{1}%【同试卷三第二(1)题】


\useproblem{3}{2}{2}%【同试卷三第二(2)题】


\useproblem{3}{2}{3}%【同试卷三第二(3)题】


\begin{problem}
设$A$, $B$, $C$三个事件两两独立，则$A$, $B$, $C$相互独立的充分必要条件是\pickout{}
\begin{abcd}
\item $A$与$BC$独立．
\item $AB$与$A \cup C$独立．
\item $AB$与$AC$独立．
\item $A \cup B$与$A \cup C$独立．
\end{abcd}
\end{problem}


\useproblem{3}{2}{5}%【同试卷三第二(5)题】


\makepart{}{本题满分6分}

\begin{problem}
已知$z = u^v$, $u = \ln \sqrt{x^2 + y^2}$, $v = \arctan \frac{y}{x}$，求$\dz$．
\end{problem}


\makepart{}{本题满分6分}

\begin{problem}
计算$I = \int_1^{+\infty} \frac{\dx}{\e^{1 + x} + \e^{3 - x}}$．
\end{problem}


\makepart{}{本题满分6分}%【同试卷三第五题】

\useproblem{3}{5}{0}


\makepart{}{本题满分7分}%【同试卷三第六题】

\useproblem{3}{6}{0}


\makepart{}{本题满分6分}

\begin{problem}
设$f(x,y) = \begin{cases}
  x^2 y, & 1 \le x \le 2,0 \le y \le x; \\
      0, & \text{其他}.
\end{cases}$求$\iint_D f(x,y)\dx\dy$，其中$D = \{(x,y)|x^2 + y^2 \ge 2x\}$
\end{problem}


\makepart{}{本题满分6分}%【同试卷一第九题】

\useproblem{1}{9}{0}


\makepart{}{本题满分8分}%【同试卷三第九题】

\useproblem{3}{9}{0}


\makepart{}{本题满分9分}

\begin{problem}
设矩阵$A = \begin{pmatrix}
   1 & -1 & 1 \\
   x &  4 & y \\
  -3 & -3 & 5
\end{pmatrix}$，已知$A$有三个线性无关的特征向量，$\lambda = 2$是$A$的二重特征值．
试求可逆矩阵$P$，使得$P^{-1} AP$为对角形矩阵．
\end{problem}


\makepart{}{本题满分8分}

\begin{problem}
设二维随机变量$(X,Y)$的密度函数为$f(x,y) = \frac{1}{2}\left[\phi_1(x,y) + \phi_2(x,y) \right]$，
其中$\phi_1(x,y)$和$\phi_2(x,y)$都是二维正态密度函数，
且它们对应的二维随机变量的相关系数分别为$\frac{1}{3}$和$- \frac{1}{3}$，
它们的边缘密度对应的随机变量的数学期望都是$0$，方差都是$1$．
\begin{enumerate}
\item 求随机变量$X$和$Y$的密度函数$f_1(x)$和$f_2(x)$．
\item 求$X$和$Y$的相关系数$\rho$．
\item 问$X$和$Y$是否独立？为什么？
\end{enumerate}
\end{problem}


\makepart{}{本题满分8分}%【同试卷三第十二题】

\useproblem{3}{12}{0}


\end{document}
